% GENERATED FILE. Edit canonical lessons, then run npm run generate:books.
% canonical-source-hash: fnv1a64:d203a293d157e9c0
% canonical-lessons: FR-C288-pickpocket, FR-C288-arnaque, FR-C288-escroc, FR-C288-escroquer, FR-C288-escroquerie

\chapter{Attention aux escrocs}
\label{ch:fr-attention-aux-escrocs}
\hlchaptermodality{288}

\begin{chapteropening}
\textbf{By the end of this chapter:} \emph{I can talk about thieves and tricksters: pickpockets, scams, con artists, swindling and fraud.}

Complete the last lesson of chapter 288: I can talk about thieves and tricksters: pickpockets, scams, con artists, swindling and fraud.
\end{chapteropening}

\section[le pickpocket]{le pickpocket — the pickpocket}

\label{lesson:FR-C288-pickpocket}



\begin{quote}
Before the new one: say the French for the forced entry, then the French for the loot.
\end{quote}

\subsection*{You'll want to know: le pickpocket}
\textbf{le pickpocket} — \textquotedblleft{}the pickpocket\textquotedblright{}.

\begin{quote}
Un pickpocket m'a pris mon téléphone dans le métro.
\end{quote}

\begin{quote}
\emph{A pickpocket took my phone on the underground.}
\end{quote}

Borrowed from English. Signs warn: \emph{attention aux pickpockets}.

\begin{cousinweb}
\emph{pickpocket} was borrowed from English \emph{pickpocket} (EtymDB), one who picks pockets.
\end{cousinweb}

\begin{grammarlens}[title={What you've built}]
One more word for telling, explaining and describing at B1.
\end{grammarlens}

\subsection*{Your turn}
\begin{itemize}
\raggedright
  \item \emph{Say it:} \emph{le pickpocket}
  \item \emph{Say it:} \emph{le pickpocket}, once more
  \item \emph{Say it:} a sentence of your own with \emph{le pickpocket}
  \item \emph{Recall:} say the French for the forced entry, then the French for the loot, then say \emph{le pickpocket} again
\end{itemize}

\begin{tcolorbox}[breakable,colback=teal!4,colframe=teal!35!black,title={Before you move on}]
What does le pickpocket mean? (\textquotedblleft{}The pickpocket\textquotedblright{}.) Say it once more.
\end{tcolorbox}
\section[l'arnaque]{l'arnaque — the scam, the rip-off (informal)}

\label{lesson:FR-C288-arnaque}



\begin{quote}
Before the new one: say the French for the loot, then the French for the pickpocket.
\end{quote}

\subsection*{You'll want to know: l'arnaque}
\textbf{l'arnaque} — \textquotedblleft{}the scam, the rip-off (informal)\textquotedblright{}.

\begin{quote}
Ce restaurant est une vraie arnaque.
\end{quote}

\begin{quote}
\emph{This restaurant is a real rip-off.}
\end{quote}

Informal. The verb is \emph{arnaquer}, to rip off.

\begin{cousinweb}
\emph{arnaque} is built on the slang verb \emph{arnaquer} \textquotedblleft{}to rip off\textquotedblright{} (EtymDB); the origin of the verb is not settled.
\end{cousinweb}

\begin{grammarlens}[title={What you've built}]
One more word for telling, explaining and describing at B1.
\end{grammarlens}

\subsection*{Your turn}
\begin{itemize}
\raggedright
  \item \emph{Say it:} \emph{l'arnaque}
  \item \emph{Say it:} \emph{l'arnaque}, once more
  \item \emph{Say it:} a sentence of your own with \emph{l'arnaque}
  \item \emph{Recall:} say the French for the loot, then the French for the pickpocket, then say \emph{l'arnaque} again
\end{itemize}

\begin{tcolorbox}[breakable,colback=teal!4,colframe=teal!35!black,title={Before you move on}]
What does l'arnaque mean? (\textquotedblleft{}The scam, the rip-off (informal)\textquotedblright{}.) Say it once more.
\end{tcolorbox}
\section[l'escroc]{l'escroc — the swindler, the con artist}

\label{lesson:FR-C288-escroc}



\begin{quote}
Before the new one: say the French for the pickpocket, then the French for the scam.
\end{quote}

\subsection*{You'll want to know: l'escroc}
\textbf{l'escroc} — \textquotedblleft{}the swindler, the con artist\textquotedblright{}.

\begin{quote}
L'escroc se faisait passer pour un policier.
\end{quote}

\begin{quote}
\emph{The con artist was passing himself off as a police officer.}
\end{quote}

The final \emph{c} is silent.

\begin{cousinweb}
\emph{escroc} was borrowed from Italian \emph{scrocco} \textquotedblleft{}sponging, living at others' cost\textquotedblright{} (EtymDB).
\end{cousinweb}

\begin{grammarlens}[title={What you've built}]
One more word for telling, explaining and describing at B1.
\end{grammarlens}

\subsection*{Your turn}
\begin{itemize}
\raggedright
  \item \emph{Say it:} \emph{l'escroc}
  \item \emph{Say it:} \emph{l'escroc}, once more
  \item \emph{Say it:} a sentence of your own with \emph{l'escroc}
  \item \emph{Recall:} say the French for the pickpocket, then the French for the scam, then say \emph{l'escroc} again
\end{itemize}

\begin{tcolorbox}[breakable,colback=teal!4,colframe=teal!35!black,title={Before you move on}]
What does l'escroc mean? (\textquotedblleft{}The swindler, the con artist\textquotedblright{}.) Say it once more.
\end{tcolorbox}
\section[escroquer]{escroquer — to swindle, to con}

\label{lesson:FR-C288-escroquer}



\begin{quote}
Before the new one: say the French for the scam, then the French for the swindler.
\end{quote}

\subsection*{You'll want to know: escroquer}
\textbf{escroquer} — \textquotedblleft{}to swindle, to con\textquotedblright{}.

\begin{quote}
Il a escroqué des personnes âgées.
\end{quote}

\begin{quote}
\emph{He swindled elderly people.}
\end{quote}

\emph{Escroquer quelqu'un de cent euros} — to con someone out of a hundred euros.

\begin{cousinweb}
\emph{escroquer} was borrowed from Italian \emph{scroccare} \textquotedblleft{}to sponge, to get by trickery\textquotedblright{} (EtymDB).
\end{cousinweb}

\begin{grammarlens}[title={What you've built}]
One more word for telling, explaining and describing at B1.
\end{grammarlens}

\subsection*{Your turn}
\begin{itemize}
\raggedright
  \item \emph{Say it:} \emph{escroquer}
  \item \emph{Say it:} \emph{escroquer}, once more
  \item \emph{Say it:} a sentence of your own with \emph{escroquer}
  \item \emph{Recall:} say the French for the scam, then the French for the swindler, then say \emph{escroquer} again
\end{itemize}

\begin{tcolorbox}[breakable,colback=teal!4,colframe=teal!35!black,title={Before you move on}]
What does escroquer mean? (\textquotedblleft{}To swindle, to con\textquotedblright{}.) Say it once more.
\end{tcolorbox}
\section[l'escroquerie]{l'escroquerie — the fraud, the swindle}

\label{lesson:FR-C288-escroquerie}



\begin{quote}
Before the new one: say the French for the swindler, then the French for to swindle.
\end{quote}

\subsection*{You'll want to know: l'escroquerie}
\textbf{l'escroquerie} — \textquotedblleft{}the fraud, the swindle\textquotedblright{}.

\begin{quote}
Il a été arrêté pour escroquerie.
\end{quote}

\begin{quote}
\emph{He was arrested for fraud.}
\end{quote}

The formal and legal word; \emph{l'arnaque} is the everyday one.

\begin{cousinweb}
\emph{escroquerie} is built on \emph{escroquer} \textquotedblleft{}to swindle\textquotedblright{} (EtymDB), which was borrowed from Italian \emph{scroccare}.
\end{cousinweb}

\begin{grammarlens}[title={What you've built}]
One more word for telling, explaining and describing at B1.
\end{grammarlens}

\subsection*{Your turn}
\begin{itemize}
\raggedright
  \item \emph{Say it:} \emph{l'escroquerie}
  \item \emph{Say it:} \emph{l'escroquerie}, once more
  \item \emph{Say it:} a sentence of your own with \emph{l'escroquerie}
  \item \emph{Recall:} say the French for the swindler, then the French for to swindle, then say \emph{l'escroquerie} again
\end{itemize}

\begin{tcolorbox}[breakable,colback=teal!4,colframe=teal!35!black,title={Before you move on}]
What does l'escroquerie mean? (\textquotedblleft{}The fraud, the swindle\textquotedblright{}.) Say it once more.
\end{tcolorbox}
